\begin{remark*}
Mathematically, a metric space is usually presented as a pair $(X,d)$, where
$X$ is a set and
\[
  d : X \times X \to \mathbb{R}
\]
is a metric satisfying the metric space axioms: nonnegativity, identity of
indiscernibles, symmetry, and the triangle inequality.

In Lean, we formalize this by separating the two pieces. First, we define a
structure \texttt{Metric Point}. Here \texttt{Point} is the Lean type playing
the role of the mathematical set $X$, and the field
\[
  \texttt{distance : Point -> Point -> Real}
\]
is the Lean version of the function $d : X \times X \to \mathbb{R}$. The metric
axioms are attached as additional fields of the same structure, so a value of
type \texttt{Metric Point} is not merely a distance function; it is a distance
function together with proofs that it satisfies the metric axioms.

Then we define a metric space as a structure
\[
  \texttt{MetricSpace Point}
\]
with a field
\[
  \texttt{metric : Metric Point}.
\]
Thus a Lean metric space is a carrier type \texttt{Point} equipped with a chosen
metric on that type. This is useful because the same underlying type may carry
different metrics. Switching the metric means changing the value of the
\texttt{metric} field while keeping the same carrier type.

The declaration
\[
  \texttt{variable \{Point : Type u\} (M : MetricSpace Point)}
\]
sets up a local context for theorem statements. It says: let \texttt{Point} be
an arbitrary type, and let \texttt{M} be an arbitrary metric space whose points
have type \texttt{Point}. The braces around \texttt{Point} make it an implicit
argument, so Lean will usually infer it from \texttt{M} or from a term such as
\texttt{x : Point}. The parentheses around \texttt{M} make the metric space an
explicit argument.

This avoids repeating the same parameters in every theorem. Instead of writing
\[
\begin{gathered}
  \texttt{theorem distance\_self}\\
  \texttt{\ \ \ \ \{Point : Type u\} (M : MetricSpace Point) (x : Point) :}\\
  \texttt{\ \ \ \ M.distance x x = 0 := by sorry,}
\end{gathered}
\]
we may first declare the variables and then write simply
\[
\begin{gathered}
  \texttt{variable \{Point : Type u\} (M : MetricSpace Point)}\\
  \texttt{theorem distance\_self (x : Point) : M.distance x x = 0 := by sorry.}
\end{gathered}
\]
Inside this context, an expression such as \texttt{M.distance x y} means the
distance from \texttt{x} to \texttt{y} according to the particular metric carried
by \texttt{M}.

The first generic theorems about \texttt{M.distance} are discharged directly
from the axiom fields stored in \texttt{M.metric}. For example, the axiom field
\[
  \texttt{M.metric.distance\_eq\_zero\_iff\_equal x y}
\]
has type
\[
  \texttt{M.metric.distance x y = 0 <-> x = y}.
\]
Since \texttt{M.distance x y} is defined to be \texttt{M.metric.distance x y},
this axiom can be used to prove the corresponding theorem stated with
\texttt{M.distance}.

Lean treats an if-and-only-if proof as having two directions. If
\[
  \texttt{h : A <-> B},
\]
then \texttt{h.mp} proves the forward direction \texttt{A -> B}, and
\texttt{h.mpr} proves the reverse direction \texttt{B -> A}. Thus, in the
theorem
\[
\begin{gathered}
  \texttt{theorem equal\_of\_distance\_eq\_zero}\\
  \texttt{\ \ \ \ \{x y : Point\} (hypothesis : M.distance x y = 0) : x = y :=}\\
  \texttt{\ \ (M.metric.distance\_eq\_zero\_iff\_equal x y).mp hypothesis,}
\end{gathered}
\]
the named assumption \texttt{hypothesis} supplies the left side of the
if-and-only-if, and \texttt{.mp} converts it into the right side. These theorems
are therefore convenience lemmas, not additional axioms: they repackage the
metric axioms using the notation attached to the particular metric space
\texttt{M}.

We call these requirements ``axioms'' because, in the mathematical definition
of a metric, they are the basic laws that a proposed distance function must
satisfy in order to count as a metric. This does not mean that we are adding new
global axioms to Lean's logical foundation. In our Lean structure, the metric
axioms are fields whose values must be supplied whenever we construct a
\texttt{Metric Point}. For a concrete example such as
\texttt{realAbsoluteValueMetric}, the fields
\texttt{distance\_nonnegative}, \texttt{distance\_eq\_zero\_iff\_equal},
\texttt{distance\_symmetric}, and \texttt{distance\_triangle} are proof
obligations. To build the example honestly, we must prove that the proposed
formula \texttt{|x - y|} satisfies each one. Thus the word ``axiom'' here means
``defining law of the mathematical structure,'' not ``unproved assumption added
to the proof system.''
\end{remark*}

\begin{example*}
The real line $\mathbb{R}$ becomes a metric space with the usual absolute-value
metric
\[
  d(x,y) = |x-y|.
\]
In Lean, the carrier type is \texttt{Real}, and the metric is defined by
\[
  \texttt{distance x y := |x - y|}.
\]
The metric-space axioms are then attached as proof fields of the structure
\texttt{Metric Real}. At this stage, those proofs may be left as \texttt{sorry}
while we focus on setting up the definitions.

The real plane $\mathbb{R}^2$ becomes a metric space with the Euclidean metric
\[
  d(p,q)
  =
  \sqrt{(p_1-q_1)^2 + (p_2-q_2)^2}.
\]
In Lean, we represent $\mathbb{R}^2$ as the product type
\[
  \texttt{Real × Real}.
\]
If \texttt{p q : Real × Real}, then \texttt{p.1} and \texttt{p.2} are the two
coordinates of \texttt{p}, and similarly for \texttt{q}. Thus the Lean distance
function is
\[
  \texttt{Real.sqrt ((p.1 - q.1) \char`^ 2 + (p.2 - q.2) \char`^ 2)}.
\]
This definition is marked \texttt{noncomputable} in Lean because the real square
root function is not executable data; it is a mathematical object supplied by
the real-number theory.
\end{example*}
